\section{回転から見直す漸化式}\label{節：回転から見直す}
\subsection{複素数の掛け算と行列}

複素数$u+iv$を$x+iy$に掛けると
\[
 (u+iv)(x+iy)=(ux-vy)+i(vx+uy).
\]
実部と虚部を縦に並べれば,
\[
 \begin{pmatrix}x'\\y'\end{pmatrix}
 =
 \begin{pmatrix}u&-v\\v&u\end{pmatrix}
 \begin{pmatrix}x\\y\end{pmatrix}.
\]
$u=\rho\cos\theta,\,v=\rho\sin\theta$なら,\,これは
\[
 \rho
 \begin{pmatrix}
 \cos\theta&-\sin\theta\\
 \sin\theta&\cos\theta
 \end{pmatrix}
\]
である. 複素数の掛け算と,\,二成分の実数の行列更新が一致した.

この行列の特性方程式は
\[
 X^2-2uX+(u^2+v^2)=0
\]
で,\,根は$u\pm iv$である.
実数の行列に非実数の固有値が現れることは,\,回転の倍率を複素数で記録したことに対応する.

ただし非実数の固有値を持つ一般の実行列が,\,元の座標でこの形に見えるとは限らない.
\sectionref{節：複素数の状態}のように,\,座標を変えて初めて回転と拡大・縮小になる場合がある.

\subsection{1次分数型も行列から生まれる}\label{節：分数と行列の合流}

1次分数型
\[
 a_{n+1}=\frac{pa_n+q}{ra_n+s}
\]
を,\,今度は二つの数の比で持ってみよう.
$(x_1,y_1)=(a_1,1)$から
\[
 \begin{pmatrix}x_{n+1}\\y_{n+1}\end{pmatrix}
 =M\begin{pmatrix}x_n\\y_n\end{pmatrix},
 \qquad
 M=\begin{pmatrix}p&q\\r&s\end{pmatrix}
\]
と更新すれば,
\[
 \frac{x_{n+1}}{y_{n+1}}
 =\cfrac{p\frac{x_n}{y_n}+q}{r\frac{x_n}{y_n}+s}.
\]
従って,\,分母が非零の範囲では$a_n=\frac{x_n}{y_n}$となる.
一次の連立更新を比として記録すると,\,1次分数型が生まれるのである.

二つの成分を同じ非零数で倍しても比は変わらない.
長さと矢印の正負を区別せず,\,原点を通る直線としての方向を記録するこの考え方を,\,\textbf{射影的な見方}という.
分母$y_n=0$は,\,成分比を有限の実数で記録できなくなる方向に当たる.
元の実数列で定義できる範囲と,\,行列で状態を追える範囲は区別する.

不動点$\alpha$については
\[
 M\begin{pmatrix}\alpha\\1\end{pmatrix}
 =(r\alpha+s)\begin{pmatrix}\alpha\\1\end{pmatrix}.
\]
従って,\,\textbf{不動点は固有値ではなく,\,固有ベクトルの成分比}である.
対応する固有値は$\lambda_\alpha=r\alpha+s$であり,\,役割が異なる.
\NewFigProjective

\subsection{差の比は固有成分の比}

異なる二つの不動点$\alpha,\beta$に対し,\,状態を
\[
 \begin{pmatrix}x_n\\y_n\end{pmatrix}
 =c_n\begin{pmatrix}\alpha\\1\end{pmatrix}
 +d_n\begin{pmatrix}\beta\\1\end{pmatrix}
\]
と分ける. 固有値を$\lambda_\alpha,\lambda_\beta$とすれば
\[
 c_{n+1}=\lambda_\alpha c_n,\qquad
 d_{n+1}=\lambda_\beta d_n.
\]
成分を比較すると
\[
 x_n-\alpha y_n=(\beta-\alpha)d_n,
\]
\[
 x_n-\beta y_n=(\alpha-\beta)c_n.
\]
従って,\,比が定義できる場合には
\[
 \frac{a_n-\alpha}{a_n-\beta}=-\frac{d_n}{c_n}
\]
となり,\,その公比は$\frac{\lambda_\beta}{\lambda_\alpha}$である.
不動点の和$r(\alpha+\beta)=p-s$から
\[
 \lambda_\beta=p-r\alpha,\qquad
 \lambda_\alpha=p-r\beta
\]
なので,\,前の1次分数型の公比と一致する.

本節でも$r\ne0,\,ps-qr\ne0$を仮定する.
このとき$M$は逆行列を持ち,\,固有値も非零である.
不動点を初項とする定数列を除けば,\,必要な成分比を取り続けられる.
元の更新の分母が零にならない条件は別に確かめる.

実係数で不動点が共役な非実数なら,\,固有値も共役である.
両者は絶対値が等しいため,\,\textbf{比を取ると共通の拡大率が消えて回転だけが残る}.
3項間漸化式では回転に$\rho^{n-1}$が伴ったが,\,1次分数型の円周上の座標ではその大きさを除いていた.

不動点が重なる場合は,\,行列も二つの固有方向に分けられない.
$r\ne0$だから$M$は単位行列の定数倍ではなく,\,固有方向は一本になる.
その寄与の蓄積を比で読み直すと,\,逆数を取った等差型が現れる.
第2章の重解と,\,第3章の重なった不動点がここでつながった.

\subsection{一つの例を三通りに記録する}

特に
\[
 M=\begin{pmatrix}1&1\\-1&1\end{pmatrix}
\]
を使おう.
状態の更新は$x_{n+1}=x_n+y_n,\,y_{n+1}=-x_n+y_n$である.
複素数を$z_n=y_n+ix_n$と選ぶと,
\[
 z_{n+1}=(1+i)z_n.
\]
この座標では,\,毎回$\sqrt2$倍されながら反時計回りに$\frac{\pi}{4}$回転する.
各成分は特性方程式$X^2-2X+2=0$に対応し,
\[
 x_{n+2}=2x_{n+1}-2x_n
\]
などを満たす.

一方,\,成分比$a_n=\frac{x_n}{y_n}$を記録すると,
\[
 a_{n+1}=\frac{a_n+1}{1-a_n}.
\]
共通の拡大率が消え,\,傾きの変化だけが残る.
さらに$w_n=\Phi(a_n)$で円周へ移すと,
\[
 w_{n+1}=iw_n.
\]
角度を2倍して記録する$\Phi$によって,\,回転角は$\frac{\pi}{2}$になる.

$(x_1,y_1)=(\frac12,1)$なら,\,成分比は先ほどの4周期の列となる.
行列では4回後の状態は$M^4\boldsymbol v_1=-4\boldsymbol v_1$であり,\,元の状態そのものには戻らない.
しかし両成分が同じ倍率を受けるため比は戻る.
\textbf{状態の周期と,\,比の周期は同じとは限らない}.

\subsection{何を記録しているか}

同じ動きを別の形で記録できることを確かめた.
一方で,\,式の見た目が似ていても更新の仕組みが異なる場合もある.

\begin{itemize}
 \item $z_{n+1}=\lambda z_n$は,\,一定角の回転と一定倍率の拡大・縮小である.
 \item $z_{n+1}=z_n^d$は,\,偏角を$d$倍し,\,絶対値も$d$乗する. 単位円上でも一定角の回転とは異なる.
 \item 実部や傾きだけを記録すると,\,異なる点が同じ値になることがある. 何の情報を残し,\,何を区別しなくなったかを確認する.
\end{itemize}

複素平面への置き換えが,\,常に全体の座標変換になるわけではない.
例えば余弦は上下二点を同じ横座標に送る.
それでも,\,元の更新と整合する点の動きを一つ作れば,\,数列の式を導ける.
一対一の座標変換と,\,一部の情報を読み取る操作を区別して使おう.

\begin{bookpoint}
漸化式の形は,\,動きそのものだけでなく,\,\textbf{その動きの何を数として記録したか}によって決まる.
不動点からのずれ,\,固有方向の成分,\,成分比,\,実部,\,傾き.
追う量を選び直すと,\,別々に見えた解法のつながりが見えてくる.
\end{bookpoint}
